\documentclass[a4paper, 11pt]{article}

\usepackage[swedish]{babel}
\usepackage[utf8]{inputenc}
\usepackage[T1]{fontenc}
\usepackage{lmodern}

\usepackage{tikz}

\tikzset{
  excl/.style = {circle, draw=black, inner sep=0pt, minimum size=4pt}
}

\tikzset{
  incl/.style = {circle, draw=black, fill=black, inner sep=0pt, minimum size=4pt}
}


\usepackage{pgfplots}
\pgfplotsset{compat=1.18}


\usepackage{graphicx}
\usepackage{caption}
\usepackage{subcaption}

\usepackage{enumitem}

\usepackage{multicol}

\usepackage{amsmath}
\usepackage{amsfonts}

\begin{document}

\title{
    \textbf{MM2001: Inlämning 5}
}
\author{Johan Montelius}
\date{HT 2026}

\maketitle


\section*{(i)}

För varje reellt tal a, bestäm antalet lösningar till följande linjära ekvationssystem:

$$
\begin{cases}
(4 + a)x_1 + 2x_2 - x_3 = 1 \\
2x_1 + (1 + a)x_2 - 2x_3 = 2 \\
-x_1 + 2x_2 + (4 + a)x_3 = 1 \\
\end{cases}
$$

Vibörjar med att lösa ut $x_3$ ur den första ekvationen:

\begin{align*}
  (4 + a)x_1 + 2x_2 - x_3 &= 1 \\  
   x_3 &= (4 + a)x_1 + 2x_2 - 1.\\  
\end{align*}

Sen substituerar vi $x_3$ i den andra ekvationen och löser ut $x_2$:

\begin{align*}
2x_1 + (1 + a)x_2 - 2x_3 &= 2 \\
2x_1 + (1 + a)x_2 - 2((4 + a)x_1 + 2x_2 - 1) &= 2\\
2x_1 + x_2 + ax_2 + 8x_1 + 2ax_1 + 4x_2 - 2 &= 2 \\
10x_1 + (5 + 3a)x_2 &= 4.\\
x_2 &= \frac{4 - 10x_1}{5 + 3a}
\end{align*}

Titt sist subsituerar vi $x_3$ och $x_2$ i den tredje ekvationen och löser ut $a$:



\begin{align*}
  -x_1 + 2x_2 + (4 + a)x_3 &= 1 \\
  -x_1 + 2\frac{4 - 10x_1}{5 + 3a} + (4+a)((4 + a)x_1 + 2x_2 - 1) &= 1\\
  -x_1 + \frac{8 - 20x_1}{5 + 3a} + (4+a)^2 x_1 + (8 + 2a)x_2 - (4+a) &= 1\\
  -x_1 + \frac{8 - 20x_1}{5 + 3a} + (4+a)^2 x_1 + (8 + 2a)(\frac{4 - 10 x_1}{5 + 3a}) - (4+a) &= 1
\end{align*}


\section*{(ii)}

{\em Hur många lösningar kan ett linjärt ekvationssystem ha? Ange även för vart och ett av
fallen ett (enkelt!) exempel på ett ekvationssystem med fyra variabler (om exemplet
inte är enkelt, behövs en motivering för vilket fall det tillhör). }

Ett ekvationssystem har: inga, en unik eller oändligt med lösningar.
Ett exemple på ett ekvationssystem utan lösningar är:

$$
\begin{cases}
x_1 + 2x_2 +3x_3 + 4x_4= 50 \\
3x_3 + 4x_4 = 37 \\
3x_3 = 3 \\
4x_4 = 4\\
\end{cases}
$$

Den tredje och fjärde ekvationen ger att $x_3$ och $x_4$ är lika med
$1$ men detta motsägs av den andra ekvationen. Ett exempel då vi har
en unik lösning är:


$$
\begin{cases}
x_1 + 2x_2 +3x_3 + 4x_4= 50\\
3x_3 + 4x_4 = 37 \\
3x_3 = 3 \\
2x_2 + 4x_4 = 30\\
\end{cases}
$$

Vi har att $x_3$ är givet (1) vilket då medför att $x_4$ är känt (9)
och det i sin tur ger värdet på $x_2$ (3). Den första ekvationen ger
slutligen ett unikt värde på $x_1$ (1).

Ett system med oändligt med lösningar måste lämna en variabel fri. Vi
kan ta samma ekvationssystem men utelämna en av de fyra ekvationerna. Utelämnar vi sen första:

$$
\begin{cases}
3x_3 + 4x_4 = 37 \\
3x_3 = 3 \\
2x_2 + 4x_4 = 30\\
\end{cases}
$$

så kan $x_1$ anta ett godtyckligt värde. Övriga variabler är dock
bundna så ett lite mer intressant exempel är att utelämna den tredje
ekvationen:

$$
\begin{cases}
x_1 + 2x_2 +3x_3 + 4x_4= 50\\
3x_3 + 4x_4 = 37 \\
2x_2 + 4x_4 = 30.\\
\end{cases}
$$

Nu är $x_3$ fri och övriaga är beroende av dess värde. 



\section*{(iii)}

{\em Hur många lösningar kan ett linjärt ekvationssystem av 5 ekvationer i 6 variabler ha?}

Inga eller oändligt - vi kan redan med två ekvationer ha en
motsägelse. Vi kan däremot inte fixera en unik lösning om vi bara har
fem ekvationer. Fem variabler kommer vara bestömda eller bestämda
givet den sjätte variabeln. Denna variabel kan väljas fritt och vi har
då oäandligt med lösningar.


\end{document}


